\subsection{Convex cones}

DEFINITION 3. --- A subset C of an affine space E is a cone with vertex \(x_0\) if C is invariant for all homotheties of centre \(x_0\) and ratio \(> 0\) .

We shall suppose in this No.~and in the one following, that we have chosen the vertex of the cone being considered, as origin in E; i.e.~we suppose that E is a vector space, and when we speak of a cone, it is to be understood that this cone has vertex 0. The set of points of the form \(\lambda a\) for \(\lambda > 0\) (resp. \(\lambda \geqslant 0\) ), where \(a\) is a non-null vector, is called an open half line (resp. closed half-line) originating at 0.

A cone C of vertex 0 is said to be pointed if \(0 \in C\) , and non-pointed otherwise. A pointed cone is either the single point \(\{0\}\) or is the union of a set of closed half-lines originating at 0. A non-pointed cone is the union (possibly empty) of open half lines originating at 0. If C is a non-pointed cone, then \(C \cup \{0\}\) is a pointed cone. If C is a pointed cone, then \(C - \{0\}\) is a non-pointed cone.

If C is a non-pointed convex cone, then \(C \cup \{0\}\) is a pointed convex cone. However, if C is a pointed convex cone, \(C - \{0\}\) is not necessarily convex. We say that a pointed convex cone is proper if it does not contain any line passing through 0. Then

PROPOSITION 9. --- A pointed convex cone C is proper if and only if the non-pointed cone C', which is the complement of 0 in C, is convex.

If C contains a line through 0 then clearly \(C'\) is not convex. Suppose now that C is proper and let x, y be two points of \(C'\) . The closed segment with end points x, y is contained in C; if it contains 0 then \(\lambda x + (1 - \lambda)y = 0\) for some \(\lambda\) with \(0 < \lambda < 1\) , therefore \(x = \mu y\) with \(\mu < 0\) . Thus C contains the line through 0 and x, contrary to hypothesis.

PROPOSITION 10. --- A subset C of E is a convex cone if and only if \(C + C \subset C\) and \(\lambda C \subset C\) for all \(\lambda > 0\) .

For the condition \(\lambda C \subset C\) for all \(\lambda > 0\) characterises the cones. If C is convex we have \(C + C = \frac{1}{2}C + \frac{1}{2}C = C\) (II, p.~8, Remark). Conversely, if the cone C is such that \(C + C \subset C\) , then for \(0 < \lambda < 1\) , we have \(\lambda C + (1 - \lambda)C = C + C \subset C\) , which shows that C is convex.

COROLLARY 1. --- If C is a non-empty convex cone, the vector space generated by C is the set C - C (the set of points x - y where x, y vary in C).

For, if V = C - C, then V is non empty, we have \(\lambda V = V\) for all \(\lambda \neq 0\) , and \(V + V = C + C - (C + C) \subset C - C = V\) , which shows that V is a vector subspace. Finally every vector subspace that contains C also contains V.

COROLLARY 2. --- If C is a pointed convex cone, the largest vector subspace contained in C is the set \(C \cap (-C)\) .

For, if \(W = C \cap (-C)\) , then W is non-empty and \(\lambda W = W\) for all \(\lambda \neq 0\) , also

\[
\mathrm{W} + \mathrm{W} \subset (\mathrm{C} + \mathrm{C}) \cap (- (\mathrm{C} + \mathrm{C})) \subset \mathrm{C} \cap (- \mathrm{C}) = \mathrm{W},
\]

which shows that W is a vector subspace. Clearly every vector subspace contained in C is also contained in W.

Obviously, if f is a linear mapping of E in a vector space F, then \(f(\mathbf{C})\) , the image of a convex cone C in E, is a convex cone in F. Every intersection of convex cones (with vertex 0) in E is a convex cone. For every subset A of E the intersection of convex cones containing A (these exist, E itself is one such cone) is the smallest convex cone that contains A; it is called the convex cone generated by A.

PROPOSITION 11. --- Let \((\mathbf{C}_{\iota})_{\iota \in \mathbf{I}}\) be a family of convex cones in E; the convex cone generated by the union of the \(\mathbf{C}_{\iota}\) is identical with the set of points \(\sum_{\iota \in \mathbf{J}} x_{\iota}\) , where \(\mathbf{J}\) is any finite

subset of I and \(x_{\iota} \in \mathbf{C}_{\iota}\) for all \(\iota \in \mathbf{J}\) .

In fact, it is obvious that C, the set of such points, is a convex cone containing the union of the \(C_{\iota}\) , and that it is contained in any convex cone which contains this union.

COROLLARY. --- For any subset A of E, the convex cone generated by A, is identical with the set of linear combinations \(\sum_{i\in J}\lambda_{i}x_{i}\) , where \((x_{i})_{i\in J}\) is any finite non-empty family of points of A, and where \(\lambda_{i}>0\) for all \(i\in J\) .

It is sufficient to see that, if a convex cone contains a point \(x \neq 0\) of A then it also contains the half-line \(C_{x}\) of the points \(\lambda x\) where \(\lambda\) varies in the set of positive numbers and that \(C_{x}\) is a convex cone.

PROPOSITION 12. --- If A is a convex set in E, then the convex cone generated by A is identical with \(C = \bigcup_{\lambda > 0} \lambda A\) .

The set C is clearly a cone; it is sufficient to show that C is convex. Let \(\lambda x, \mu y\) be two points of C ( \(\lambda > 0, \mu > 0, x \in A, y \in A\) ). Let \(\alpha, \beta\) be two numbers \textgreater{} 0 such that \(\alpha + \beta = 1\) . Then \(\alpha \lambda x + \beta \mu y = (\alpha \lambda + \beta \mu) z\) , with \(z \in A\) , and \(\alpha \lambda + \beta \mu > 0\) ; hence \(\alpha \lambda x + \beta \mu y \in C\) .

Remarks. --- 1) With the hypotheses of prop. 12, if \(0 \notin \mathbf{A}\) , then the cone \(C\) is non-pointed, thus \(C \cup \{0\}\) is proper.

\begin{enumerate}
\def\labelenumi{\arabic{enumi})}
\setcounter{enumi}{1}
\tightlist
\item
  Let A be any convex set in E; consider the convex set \(A_{1} = A \times \{1\}\) in the space \(F = E \times R\) and the convex cone C with vertex 0 that is generated by \(A_{1}\) . Prop. 12 shows that \(A_{1}\) is the intersection of C and of the hyperplane \(E \times \{1\}\) in F. Every convex set in E can, therefore, be considered as the projection on E of the intersection of a convex cone with vertex 0 in F and the hyperplane \(E \times \{1\}\) .
\end{enumerate}

